\section{Analiza rynku}

Na rynku istnieje wiele serwisów internetowych, na których użytkownicy mogą licytować i kupować przedmioty w czasie rzecywistym.  Jak okazuje się, nawet na największych platformach, nie ma możliwości podglądu obiektu w trzech wymiarach.
Witryny takie jak allegro.pl, amazon.pl czy Erli posiadają prostą funkcjonalność podglądu zdjęć, nie wyróżniając się innowacyjnym rozwiązaniem jakim jest rendering 3D licytowanych obiektów.
Obrazy 2D zapewniają statyczny widok i nie oddają skomplikowanych szczegółów, kątów lub faktur. Prowadzi to do tego, że nie znamy rzeczywistego stanu przedmiotu, przez co potencjalny konsument, może zostać wprowadzony w błąd, podczas zakupu. 
Technologia ta pozwala nam na dokładną inspekcję wybranego przez nas przedmiotu, poprzez obracanie obiektem, możliwością przybliżania oraz oddalania się od obiektu i renderingu w bardzo wysokiej jakości.

\subsection{allegro.pl}
Jest to największy serwis aukcyjny w Polsce, założony w 1999 roku przez Surf Stop Shop. Ponad 33 \% konsumentów zaczyna właśnie wyszukiwanie produktów na wyżej wymienionym serwisie. Podgląd obiektów na tej stronie jest bardzo prosty, gdyz pozwala nam tylko na podgląd zdjęć, bez możliwosci zaawansowanego podglądu licytowanego obiektu. Interfejs użytkownika jest prosty i bardzo intuicyjny, a funkcjonalność licytowania pozwala wielu użytkownikom licytować te same przedmioty bez żadnych problemów. Na tym portalu każdy użytkownik może zostać sprzedawcą i wystawić własne przedmioty, nie korzystając z żadnego pośrednika, co ułatwia kontakt pomiędzy konsumentem a sprzedawcą, dzięki czemu każdy problem dotyczący sprzedaży może zostać rozwiązany w prosty sposób.

%wsadzic tu focie

\subsection{amazon.pl}
Amazon to jeden z największych serwisów e-commerce na całym świecie. Startował jako księgrania internetowa, której asortyment systematycznie powiększał się, aż do dzisiejszych rozmiarów. Jedyną funkcjonalnością podglądu kupowanego przedmiotu jest przbyliżenie obrazka, który zapewnił nam sprzedający. Nie mamy możliwości podglądu realnego stanu licytowanego przedmiotu. Sam proces kupowania czy licytowania jest bardzo intuicyjny i szybki, przez co użytkownik nie ma problemów. Amazon ma bardzo dobrze rozbudowaną obsługę kilenta, przez co wszystkie problemy są rozwiązywane szybko i bezproblemowo. 

%wsadzic tu focie

\subsection{erli.pl}
Założona w 2020 polska firma, która zajmuje się sektorem e-commerce. Zajmuje się sprzedaża przemiotów z kategorii: meble, budowa i remont, ogród, oświetlenie, narzędzia i wyposażenie domu.
Strona ma prosty i przejrzysty układ strony kupowanego przedmiotu, lecz podgląd kupowanego przedmiotu ogranicza się tylko i wyłącznie do powiększenia zdjęć oferowanych przez sprzedawce. Takie rozwiązanie może nie oddawać dobrze faktur i tekstur przedmiotu, przez co użytkownik moze zostać wprowadzony w skonfundowanie podczas zakupu. Wyświetlane zdjęcia są na małej częsci strony, co utrudnia ogląd licytowanego przedmiotu dla użytkowników, gdyz muszą powiększać stronę, aby dobrze obejrzeć zdjęcia zawarte przez sprzedającego. 

%wsadzic tu focie

